\section{Swap: a Cycle Broken by Reads at Fresh Positions}\label{sec:swap}

\subsection{The program}
\begin{lstlisting}
method swap(a: array<int, int>, b: array<int, int>, n: int)
  returns (a_: array<int, int>, b_: array<int, int>)
  requires 0 <= n
  ensures forall k: int :: 0 <= k < n ==> a_[k] == b[k] && b_[k] == a[k]
{
  a_ := a; b_ := b; i := 0;
  while i < n { t := a_[i]; a_[i] := b_[i]; b_[i] := t; i := i + 1; }
}
\end{lstlisting}
The exchange of the first $n$ elements of two arrays, element by element (the example \texttt{swap.dfy} of
\texttt{examples/dafny/state}). Each array is written with values read from the other, so neither can be
described without the other: the two arrays form a cycle of the read relation. A proof by an invariant states what
both arrays hold below and above $i$; here, no invariant is given.

\subsection{The system}
Symbolic execution substitutes $t := a_\_[i]$ into the store of $b_\_$, so the update of $a_\_$ reads $b_\_$ and the
update of $b_\_$ reads $a_\_$. The two arrays are one component and get one trace $S$ of two accumulators
(Section~\ref{sec:teq}); the index and the temporary $t$ are components of their own:
\begin{gather}
  K = \cnt(0, 1, n') \qquad I = \tr{0}{\bl{K}}{i, k \mapsto i + 1} \label{eq:sw0}\\
  S = \mathsf{trace}(a, b,\ \mathsf{for}\ \bl{K}, \bl{I}\ \mathsf{do}\ a_\_, b_\_, k, i \mapsto
      a_\_[i := b_\_[i]],\ b_\_[i := a_\_[i]]) \label{eq:sw1}\\
  T = \tr{t_0}{\bl{K}, \bl{S}.0, \bl{I}}{t, k, a_\_, i \mapsto a_\_[i]} \label{eq:sw2}\\
  \al{\bl{K}, \ldots, \bl{I}}{\ldots, i \mapsto i < n} \qquad \neg(\cur(I) < n) \label{eq:sw3}
\end{gather}
with the results $\cur(S).0$ and $\cur(S).1$. The run is \texttt{-teq -prove}; it takes 2\,s and 15 lemmas of
\texttt{teq.dfy}, no optional module. $T$ is not read by the goal and plays no part.

\subsection{What the proof needs}
\emph{The counter.} $I$ is $\cnt(0, 1, n')$ by \lemma{counter\_def}, so the index in iteration $k$ is $k$
(\lemma{earlier\_counter}), and $\cur(I) = \intf(n')$.

\emph{A view of each component as stores.} No lemma is written for a trace of two arrays. The engine reads
$\cur(S).1$ through its view (Section~\ref{sec:engine:views}): the trace of $b_\_$ alone, with the history of the
other accumulator as a further dependency,
\begin{align*}
  &\tr{b}{\bl{I}, \bl{S}.0}{b_\_, k, a_\_ \mapsto b_\_[k := a_\_[k]]}\\
  &\quad \leadsto\ \mathsf{stores}(b, \bl{I}, \mapc{\bl{I}, \bl{S}.0}{k, a_\_ \mapsto a_\_[k]}),
\end{align*}
by the definition of $\mathsf{stores}$, which the engine folds: a trace whose update stores at the keys is the
sequence of stores at these keys, of the values listed. The values still read the other array, so the cycle is
not yet broken; but the keys are now known to be the counter $0, \ldots, n' - 1$.

\emph{A read at a fresh position.} The value stored in iteration $k$ is $a_\_[k]$, the other array as it was
before iteration $k$. The view of $a_\_$ is a sequence of stores as well, and \lemma{earlier\_stores} reads an
earlier state of stores at a position:
\begin{align*}
  \mathsf{earlier}(j, \bl{\mathsf{stores}(a, ks, vs)})[k] = {}
  &\mathsf{if}\ \mathsf{occurs}(k, \mathsf{prefix}(j, ks))\\
  &\mathsf{then}\ \mathsf{earlier}(\mathsf{indexof}(k, \mathsf{prefix}(j, ks)), vs)\ \mathsf{else}\ a[k].
\end{align*}
In iteration $k$, the keys written so far are $\mathsf{prefix}(k, \bl{I}) = \bl{\cnt(0, 1, k)}$
(\lemma{prefix\_counter}), and $k$ does not occur among $0, \ldots, k - 1$ (\lemma{occurs\_counter\_up}): the
position has not been written yet, and the read is the initial value, $a_\_[k] = a[k]$. The stored values become
$\mapc{\bl{I}}{k \mapsto a[k]}$, which folds to $\mathsf{selects}(a, \bl{I})$, and symmetrically for the other
component:
\begin{gather*}
  \cur(S).0 = \cur(\mathsf{stores}(a, \bl{I}, \mathsf{selects}(b, \bl{I})))\\
  \cur(S).1 = \cur(\mathsf{stores}(b, \bl{I}, \mathsf{selects}(a, \bl{I}))).
\end{gather*}
The cycle is gone: each array is the other's initial values stored at $0, \ldots, n' - 1$.

\emph{The postcondition.} The iteration constraint gives $\intf(n') \le n$ for $0 \le n$
(\lemma{invs\_counter\_upper\_nonnegative}), and the exit $n \le \cur(I) = \intf(n')$, so $\intf(n') = n$. A read of
the final array at $0 \le k < n$ is \lemma{current\_stores}: $k$ occurs among the keys, at index $k$
(\lemma{indexof\_counter\_up}), and the value is the $k$-th of $\mathsf{selects}(b, \bl{I})$, which is $b[k]$
(\lemma{earlier\_selects}). Both conjuncts of the postcondition hold, and the goal is dropped.

\subsection{What it shows}
The cycle is resolved by the order of the writes, not by an invariant: the loop writes each position once, in
increasing order, and reads the other array only at the position it is about to write. The lemma that captures
this, \lemma{earlier\_stores}, is about a single trace of stores, and the counter lemmas decide that the position is
fresh. Nothing in the database mentions two arrays or a swap: the views turn the component of the cycle into a
trace of one accumulator, to which the lemmas of single loops apply, and the cycle dissolves once both views are
stores of known values. A loop that read the other array at a position already written, e.g.\ at $i - 1$, would
make the read an earlier stored value, a recurrence between the two arrays, which these lemmas do not close.
